%% toc: P3 26 26.4
\section{Proofs: finite signatures and graph types}
\label{sec:annex-p03-w03-p03-015-finite-signature}

Throughout, \(S_n(T)\) is the set of consistent complete \(n\)-types of
\cref{def:p03-rn-standing}, and patterns, pattern formulas and
\(\operatorname{At}_n\) are those of \cref{def:p03-w03-p03-015-patterns}.
Automorphisms preserve the truth of every formula at every tuple
(\cref{lem:iso-invariance}).

\begin{annexitem}{annex:p03-w03-p03-015-app}{thm:p03-w03-p03-015-finite-signature-application}{Patterns decide types}
\begin{proof}
\emph{Finiteness of the atomic list.} Since \(L\) has no function
symbols, every term in \(\bar x\) is a variable \(x_i\) or a constant
\(c_j\); there are \(n+k\) of them. An atomic formula is an equation
between two such terms or a relation symbol applied to
\(\operatorname{ar}(R)\) of them. As \(L\) has finitely many relation
symbols, \(\operatorname{At}_n\) is finite, with the size displayed in
\cref{def:p03-w03-p03-015-patterns}, and there are at most
\(2^{|\operatorname{At}_n|}\) patterns.

\emph{(1) Equal patterns give an automorphism.} Let \(\bar a,\bar b\in M^n\)
have the same pattern \(\Delta\). Put \(X=\{a_0,\ldots,a_{n-1}\}\) and
\(Y=\{b_0,\ldots,b_{n-1}\}\). The set
\(X\cup\{c_1^{\mathcal M},\ldots,c_k^{\mathcal M}\}\) contains \(X\) and
the constants and is closed under the (absent) functions, so it is
\(\langle X\rangle^{\mathcal M}\); it is finite, and each of its elements
is the value \(t^{\mathcal M}[\bar a]\) of some term \(t\) in \(\bar x\).
The same holds for \(\langle Y\rangle^{\mathcal M}\) and \(\bar b\).
Define \(f(t^{\mathcal M}[\bar a])=t^{\mathcal M}[\bar b]\).
For terms \(t,t'\),
\[
 t^{\mathcal M}[\bar a]=t'^{\mathcal M}[\bar a]
 \iff(t=t')\in\Delta
 \iff t^{\mathcal M}[\bar b]=t'^{\mathcal M}[\bar b],
\]
so \(f\) is well defined and injective; it is onto
\(\langle Y\rangle^{\mathcal M}\) because every element there is a term
value at \(\bar b\). Taking \(t=c_j\) shows that \(f\) fixes every
constant. For a relation symbol \(R\) of arity \(m\) and terms
\(t_1,\ldots,t_m\),
\[
 R^{\mathcal M}(t_1[\bar a],\ldots,t_m[\bar a])
 \iff R(t_1,\ldots,t_m)\in\Delta
 \iff R^{\mathcal M}(t_1[\bar b],\ldots,t_m[\bar b]).
\]
Every tuple from \(\langle X\rangle^{\mathcal M}\) has this form, so \(f\)
is an isomorphism of finitely generated substructures with
\(f(a_i)=b_i\). Ultrahomogeneity gives \(\sigma\in\operatorname{Aut}(\mathcal M)\)
extending \(f\), and \(\sigma\bar a=\bar b\). Consequently \(\bar a\) and
\(\bar b\) satisfy the same formulas in \(\mathcal M\).

\emph{(2) Quantifier elimination.} Fix \(\varphi(\bar x)\) and let
\(D_\varphi\) be the set of \(\Delta\in P_n(\mathcal M)\) such that
\(\mathcal M\models\varphi(\bar a)\) for some \(\bar a\) with pattern
\(\Delta\). By (1), \(\varphi\) then holds at every tuple with that
pattern. Let \(\bar a\in M^n\). If \(\mathcal M\models\varphi(\bar a)\),
its pattern \(\Delta^{\mathcal M}(\bar a)\) lies in \(D_\varphi\) and
\(\bar a\) satisfies \(\pi_{\Delta^{\mathcal M}(\bar a)}\). Conversely, if
\(\bar a\) satisfies \(\pi_\Delta\) with \(\Delta\in D_\varphi\), its
pattern is \(\Delta\), so \(\varphi\) holds at \(\bar a\). Hence
\[
 \mathcal M\models\forall\bar x\Bigl(\varphi(\bar x)\leftrightarrow
       \bigvee_{\Delta\in D_\varphi}\pi_\Delta(\bar x)\Bigr),
\]
and this sentence belongs to \(T\). The right side is quantifier-free.
For \(n=0\) the same argument shows that every sentence is equivalent
modulo \(T\) to the single realized pattern formula or to falsity.

\emph{(3) Types.} For \(\Delta\in P_n(\mathcal M)\) let
\(p_\Delta=\{\varphi:\Delta\in D_\varphi\}\). If \(\bar a\) has pattern
\(\Delta\), then \(p_\Delta=\operatorname{tp}^{\mathcal M}(\bar a)\): a
formula true at \(\bar a\) puts \(\Delta\) into \(D_\varphi\), and a
formula with \(\Delta\in D_\varphi\) is true at \(\bar a\) by (1).
So \(p_\Delta\in S_n(T)\). It contains \(\pi_\Delta\), and for each
\(\varphi\in p_\Delta\) the formula \(\pi_\Delta\) is one of the disjuncts
in (2), so \(T\models\forall\bar x(\pi_\Delta\to\varphi)\). Since
\(\pi_\Delta\) is consistent with \(T\), it isolates \(p_\Delta\).
If \(\Delta\neq\Delta'\), a realization of \(\Delta'\) fails
\(\pi_\Delta\), so \(\neg\pi_\Delta\in p_{\Delta'}\) and
\(p_\Delta\neq p_{\Delta'}\).

For surjectivity, let \(q\in S_n(T)\). Every tuple of \(M^n\) has a
pattern in the finite set \(P_n(\mathcal M)\), so
\(\forall\bar x\bigvee_{\Delta\in P_n(\mathcal M)}\pi_\Delta\) is true in
\(\mathcal M\) and belongs to \(T\). If \(q\) contained \(\neg\pi_\Delta\)
for every \(\Delta\in P_n(\mathcal M)\), a realization of \(q\) in a model
of \(T\) would satisfy that disjunction and all of these negations, which
is impossible. As \(q\) is complete, it contains some \(\pi_\Delta\). For
\(\varphi\in p_\Delta\), \(T\models\forall\bar x(\pi_\Delta\to\varphi)\),
so consistency of \(q\) excludes \(\neg\varphi\) and completeness gives
\(\varphi\in q\). Thus \(p_\Delta\subseteq q\). If some \(\psi\in q\)
were not in \(p_\Delta\), then \(\neg\psi\in p_\Delta\subseteq q\),
contradicting consistency. So \(q=p_\Delta\), and
\(|S_n(T)|=|P_n(\mathcal M)|\).

\emph{(4) Categoricity.} The signature \(L\) is finite, hence countable.
\(T\) is complete and satisfiable and has the infinite model
\(\mathcal M\), so \cref{def:p03-rn-standing} holds and every model of
\(T\) is infinite. By (3), \(S_n(T)\) is finite for every \(n\), and
\cref{thm:p03-rn-characterization} gives \(\omega\)-categoricity. A
countable model of \(T\) is therefore countably infinite and isomorphic
to \(\mathcal M\).

\emph{Fra\"iss\'e limits.} An infinite Fra\"iss\'e limit is, by
\cref{def:p03-w03-p03-015-ultrahomogeneous}, a countably infinite
ultrahomogeneous structure, so the four conclusions apply to it.
\end{proof}
\end{annexitem}

\begin{annexitem}{annex:p03-w03-p03-015-fn}{thm:p03-w03-p03-015-function-boundary}{Function symbols}
\begin{proof}
The domain \(\mathbb Z\) is countably infinite. The signature has no
constants, so \(\langle\varnothing\rangle^{\mathcal Z}=\varnothing\). A
nonempty subset of \(\mathbb Z\) closed under \(s\) and \(t\) contains
\(a+j\) for every integer \(j\) whenever it contains \(a\), so
\(\langle X\rangle^{\mathcal Z}=\mathbb Z\) for every nonempty \(X\).
Let \(f:\langle X\rangle^{\mathcal Z}\to\langle Y\rangle^{\mathcal Z}\)
be an isomorphism with \(X,Y\) finite. Because \(f\) is a bijection,
either both generated sets are empty or both are \(\mathbb Z\). In the
first case \(f\) is the empty map, which the identity extends. In the
second, \(f\) is a bijection of \(\mathbb Z\) commuting with \(s\) and
\(t\), that is, an automorphism of \(\mathcal Z\), and extends itself.
So \(\mathcal Z\) is ultrahomogeneous.

For \(j\geq0\), the pair \((0,j)\) satisfies \(y=s^j(x)\), and
\((0,i)\) does not for \(i\neq j\). Hence the types of the pairs
\((0,j)\) are pairwise distinct, and \(S_2(\operatorname{Th}(\mathcal Z))\)
is infinite, as also recorded in \cref{thm:p03-rn-function-boundary}.
The signature is countable, the theory is complete and has the infinite
model \(\mathcal Z\), so by \cref{thm:p03-rn-characterization} it is not
\(\omega\)-categorical.
\end{proof}
\end{annexitem}

\begin{annexitem}{annex:p03-w03-p03-015-gr}{thm:p03-w03-p03-015-graph-realization}{Class by class}
\begin{proof}
Let \((\varepsilon,\eta)\in\Gamma_n\) and list its classes
\(C_1,\ldots,C_m\) in order of their least elements. Let
\(\mathcal G\models T_{\rm RG}\). We choose distinct vertices
\(v_1,\ldots,v_m\) of \(\mathcal G\) such that
\(v_iEv_j\Leftrightarrow C_i\,\eta\,C_j\) for \(i\neq j\).

Suppose \(v_1,\ldots,v_j\) (\(0\leq j<m\)) are distinct and satisfy this
for indices up to \(j\). Let \(U\) be the list of those \(v_i\), \(i\leq j\),
with \(C_i\,\eta\,C_{j+1}\), and \(V\) the list of the remaining
\(v_i\), \(i\leq j\); put \(r=|U|\) and \(s=|V|\). The \(r+s\) vertices
in \(U\) followed by \(V\) are distinct, so the extension axiom for
\((r,s)\) gives \(z\) distinct from all of them, adjacent to each vertex
in \(U\) and to none in \(V\). (For \(j=0\) this is the axiom with
\(r=s=0\), which supplies some vertex.) Put \(v_{j+1}=z\). Adjacency
between \(v_{j+1}\) and earlier vertices is as required; by symmetry of
\(E\) it is the same in both orders.

Define \(a_i=v_\ell\) when \(i\in C_\ell\). Then \(a_i=a_j\) exactly when
\(i\,\varepsilon\,j\), because the \(v_\ell\) are distinct, and for
\([i]\neq[j]\), \(E(a_i,a_j)\) holds exactly when \([i]\,\eta\,[j]\).
So \(\bar a\) satisfies \(\gamma_{\varepsilon,\eta}\). For \(n=0\) the
empty tuple satisfies \(\top\).
\end{proof}
\end{annexitem}

\begin{annexitem}{annex:p03-w03-p03-015-gp}{thm:p03-w03-p03-015-graph-partial}{One-point extensions}
\begin{proof}
\emph{One-point extension.} Let \(f\) be a finite partial isomorphism from
\(\mathcal G\) to \(\mathcal H\) and \(b\in G\) outside its domain
\(D\). Let \(U=\{u\in D:uEb\}\) and \(W=D\setminus U\). The images
\(f(U)\) and \(f(W)\) are disjoint lists of \(|D|\) distinct vertices of
\(\mathcal H\), since \(f\) is injective. The extension axiom in
\(\mathcal H\) for \((|U|,|W|)\) gives \(z\) outside \(f(D)\), adjacent to
every vertex of \(f(U)\) and to none of \(f(W)\). Then
\(f'=f\cup\{(b,z)\}\) is injective, and for \(u\in D\),
\(uEb\Leftrightarrow u\in U\Leftrightarrow f(u)Ez\); by symmetry the same
holds for \(bEu\), and \(bEb\), \(zEz\) both fail by irreflexivity. So
\(f'\) is a finite partial isomorphism. The inverse of a finite partial
isomorphism from \(\mathcal G\) to \(\mathcal H\) is one from
\(\mathcal H\) to \(\mathcal G\), and the extension axioms hold in both
structures, so the same extension is available in the other direction.

\emph{Induction.} Formulas are built from atomic formulas by \(\neg\),
\(\land\) and \(\exists\); the other connectives are abbreviations. We prove
the equivalence for all pairs of models of \(T_{\rm RG}\), all finite
partial isomorphisms between them and all tuples from their domains, by
induction on the construction of \(\varphi\).
For \(x_i=x_j\): \(a_i=a_j\) iff \(f(a_i)=f(a_j)\), as \(f\) is injective.
For \(E(x_i,x_j)\): this is the defining property of \(f\).
Negation and conjunction follow immediately from the induction hypothesis.

Let \(\varphi=\exists y\,\psi(\bar x,y)\) and suppose
\(\mathcal G\models\psi(\bar a,b)\). If \(b\) is in the domain of \(f\),
the induction hypothesis for \(\psi\) and \(f\) gives
\(\mathcal H\models\psi(f\bar a,f(b))\). Otherwise extend \(f\) to \(f'\)
with \(f'(b)=z\) as above; the hypothesis for \(\psi\), \(f'\) and the
tuple \((\bar a,b)\) gives \(\mathcal H\models\psi(f\bar a,z)\). Either
way \(\mathcal H\models\varphi(f\bar a)\). The converse is the same
argument applied to \(f^{-1}\), with the roles of \(\mathcal G\) and
\(\mathcal H\) exchanged, which the induction hypothesis covers because it
is stated for all pairs of models.
\end{proof}
\end{annexitem}

\begin{annexitem}{annex:p03-w03-p03-015-gt}{thm:p03-w03-p03-015-graph-types}{Graph patterns as types}
\begin{proof}
\emph{(1)} Let \(\bar a\) in \(\mathcal G\) and \(\bar b\) in
\(\mathcal H\) satisfy the same \(\gamma_{\varepsilon,\eta}\). The rule
\(a_i\mapsto b_i\) is well defined and injective, because
\(a_i=a_j\), \(i\,\varepsilon\,j\) and \(b_i=b_j\) are equivalent. For
\(a_i\neq a_j\) the classes differ and both \(E(a_i,a_j)\) and
\(E(b_i,b_j)\) are decided by \([i]\,\eta\,[j]\); for \(a_i=a_j\) both
fail by irreflexivity. So the rule is a finite partial isomorphism, and
\cref{thm:p03-w03-p03-015-graph-partial} gives agreement on every formula.
The empty map is a finite partial isomorphism between any two models, so
any two models satisfy the same sentences. \(T_{\rm RG}\) has the model
\(\mathcal R\) (\cref{thm:p03-fraisse54-graph}); hence it is complete.

\emph{(2)} Fix \(\varphi(\bar x)\) and let \(D_\varphi\) be the set of
\((\varepsilon,\eta)\in\Gamma_n\) at whose realizations in \(\mathcal R\)
the formula \(\varphi\) holds; by (1) and
\cref{thm:p03-w03-p03-015-graph-realization} this does not depend on the
realization or on the model. In a model of \(T_{\rm RG}\), each tuple
satisfies exactly one \(\gamma_{\varepsilon,\eta}\)
(\cref{def:p03-w03-p03-015-graph-patterns}); it satisfies \(\varphi\)
exactly when that pattern lies in \(D_\varphi\). Therefore
\(T_{\rm RG}\models\forall\bar x(\varphi\leftrightarrow
\bigvee_{(\varepsilon,\eta)\in D_\varphi}\gamma_{\varepsilon,\eta})\).

\emph{(3)} The type of any realization of \(\gamma_{\varepsilon,\eta}\)
is, by (1), the same set \(p_{\varepsilon,\eta}\) in every model, and it
contains \(\gamma_{\varepsilon,\eta}\). For each
\(\varphi\in p_{\varepsilon,\eta}\), the pattern lies in \(D_\varphi\), so
\(T_{\rm RG}\models\forall\bar x(\gamma_{\varepsilon,\eta}\to\varphi)\);
as \(\gamma_{\varepsilon,\eta}\) is consistent with \(T_{\rm RG}\) by
\cref{thm:p03-w03-p03-015-graph-realization}, it isolates
\(p_{\varepsilon,\eta}\). Distinct patterns are jointly unsatisfiable, so
their types differ. If \(q\in S_n(T_{\rm RG})\), then since
\(T_{\rm RG}\models\forall\bar x\bigvee_{\Gamma_n}\gamma_{\varepsilon,\eta}\)
and \(\Gamma_n\) is finite, \(q\) contains some \(\gamma_{\varepsilon,\eta}\);
then \(p_{\varepsilon,\eta}\subseteq q\) because each of its formulas is
implied by \(\gamma_{\varepsilon,\eta}\) modulo \(T_{\rm RG}\), and the
two complete types coincide.

\emph{(4)} A graph pattern consists of an equivalence relation with some
number \(k\) of classes, chosen in \({n\brace k}\) ways, and a simple
graph on those \(k\) distinguishable classes, chosen in
\(2^{k(k-1)/2}\) ways, one for each subset of the \(k(k-1)/2\) unordered
pairs. Summing over \(k\) gives the formula. The values are
\(1\) for \(n=0\) and \(n=1\); \({2\brace1}+{2\brace2}\cdot2=3\);
\(1+3\cdot2+1\cdot8=15\); and \(1+7\cdot2+6\cdot8+1\cdot64=127\),
using \({4\brace2}=7\) and \({4\brace3}=6\).

\emph{(5)} Automorphisms of \(\mathcal R\) preserve every formula, so
tuples in one orbit satisfy the same \(\gamma_{\varepsilon,\eta}\).
Conversely, two tuples of \(\mathcal R\) satisfying the same
\(\gamma_{\varepsilon,\eta}\) define, as in (1), an isomorphism between
the finite induced subgraphs on their entries; ultrahomogeneity of
\(\mathcal R\) (\cref{thm:p03-fraisse54-graph}) extends it to an
automorphism carrying one tuple to the other. Every pattern is realized in
\(\mathcal R\) by \cref{thm:p03-w03-p03-015-graph-realization}, so there
are exactly \(|\Gamma_n|\) orbits.
\end{proof}
\end{annexitem}

\begin{annexitem}{annex:p03-w03-p03-015-rg}{thm:p03-w03-p03-015-random-graph-instance}{Matching patterns}
\begin{proof}
\(\mathcal R\) is countably infinite and ultrahomogeneous
(\cref{thm:p03-fraisse54-graph}), and \(L_E\) is finite with no function
or constant symbols, so \cref{thm:p03-w03-p03-015-finite-signature-application}
applies. Here \(\operatorname{At}_n\) consists of the \(n^2\) equations
\(x_i=x_j\) and the \(n^2\) atoms \(E(x_i,x_j)\).

For \(\Delta\in P_n(\mathcal R)\), realized by \(\bar a\), let
\(i\,\varepsilon_\Delta\,j\) mean \((x_i=x_j)\in\Delta\); this is the
relation \(a_i=a_j\), an equivalence relation. Let
\([i]\,\eta_\Delta\,[j]\) mean that \([i]\neq[j]\) and
\(E(x_i,x_j)\in\Delta\). It does not depend on the representatives,
because equal entries have the same neighbours; it is irreflexive by
definition and symmetric because \(E\) is. So
\((\varepsilon_\Delta,\eta_\Delta)\in\Gamma_n\), and it is the graph
pattern of \(\bar a\).

The map is injective: \(\Delta\) is recovered from
\((\varepsilon_\Delta,\eta_\Delta)\), since \((x_i=x_j)\in\Delta\) iff
\(i\,\varepsilon_\Delta\,j\), and \(E(x_i,x_j)\in\Delta\) iff
\([i]\neq[j]\) and \([i]\,\eta_\Delta\,[j]\), the case \([i]=[j]\) being
excluded by irreflexivity. It is surjective because every graph pattern
is realized in \(\mathcal R\) (\cref{thm:p03-w03-p03-015-graph-realization}),
and the pattern of a realization maps to it.

The formula \(\pi_\Delta\) contains every literal of
\(\gamma_{\varepsilon_\Delta,\eta_\Delta}\); its remaining literals are
\(\neg E(x_i,x_j)\) for \(i\,\varepsilon_\Delta\,j\), which follow from
the equalities and irreflexivity, and edge literals for pairs of indices
from distinct classes, which follow from those of
\(\gamma_{\varepsilon_\Delta,\eta_\Delta}\) by the equalities and
symmetry. So the two formulas are equivalent modulo \(T_{\rm RG}\).
\cref{thm:p03-w03-p03-015-finite-signature-application} now gives
\(|S_n(\operatorname{Th}(\mathcal R))|=|P_n(\mathcal R)|=|\Gamma_n|\),
isolation by the \(\pi_\Delta\), and \(\omega\)-categoricity.
Since \(T_{\rm RG}\) is complete (\cref{thm:p03-w03-p03-015-graph-types})
and true in \(\mathcal R\), its consequences are exactly
\(\operatorname{Th}(\mathcal R)\).
\end{proof}
\end{annexitem}

\begin{annexitem}{annex:p03-w03-p03-015-nv}{thm:p03-w03-p03-015-named-vertex}{The named vertex}
\begin{proof}
\emph{Ultrahomogeneity.} In \((\mathcal R,v)\), every generated set is
\(\langle X\rangle=X\cup\{v\}\), finite for finite \(X\). Let
\(f:X\cup\{v\}\to Y\cup\{v\}\) be an isomorphism. It fixes \(v\) and
preserves adjacency and non-adjacency, so it is an isomorphism between
finite induced subgraphs of \(\mathcal R\). By ultrahomogeneity of
\(\mathcal R\) it extends to \(\sigma\in\operatorname{Aut}(\mathcal R)\),
and \(\sigma(v)=f(v)=v\), so \(\sigma\) is an automorphism of
\((\mathcal R,v)\). The structure is countably infinite and the signature
\(\{E,c\}\) is finite with no function symbols, so
\cref{thm:p03-w03-p03-015-finite-signature-application} applies.

\emph{Count.} The terms on \(n\) slots are \(x_0,\ldots,x_{n-1},c\).
Replacing \(c\) by a new variable \(x_n\) is a bijection between
\(\operatorname{At}_n\) for \(\{E,c\}\) and \(\operatorname{At}_{n+1}\)
for \(L_E\), and the pattern of \(\bar a\) in \((\mathcal R,v)\)
corresponds to the pattern of \((\bar a,v)\) in \(\mathcal R\). So
\(P_n(\mathcal R,v)\) corresponds to the patterns of \((n+1)\)-tuples of
\(\mathcal R\) whose last entry is \(v\). Every pattern in
\(P_{n+1}(\mathcal R)\) is of this kind: if \((\bar b,w)\) realizes it,
the one-point map \(w\mapsto v\) is an isomorphism of one-vertex induced
subgraphs (both loopless), and its extension
\(\tau\in\operatorname{Aut}(\mathcal R)\) carries \((\bar b,w)\) to
\((\tau\bar b,v)\), which has the same pattern. Hence
\(|P_n(\mathcal R,v)|=|P_{n+1}(\mathcal R)|=|\Gamma_{n+1}|\) by
\cref{thm:p03-w03-p03-015-random-graph-instance}, and part (3) of the
application theorem gives \(|S_n(\operatorname{Th}(\mathcal R,v))|=|\Gamma_{n+1}|\).
The values \(1,3,15,127\) are those of
\cref{thm:p03-w03-p03-015-graph-types} for \(n+1=1,2,3,4\).
Part (4) of the application theorem gives \(\omega\)-categoricity.
\end{proof}
\end{annexitem}

\begin{annexitem}{annex:p03-w03-p03-015-effective-elimination}{thm:p03-w03-p03-015-effective-elimination}{Computing patterns and realizing witnesses}
\begin{proof}
\emph{The extension property needed for projection.}
Fix a tuple \(\bar a\) in a model \(\mathcal G\models T_{\rm RG}\),
and write \(\rho\) for its graph pattern. Let
\(\rho'\in\Gamma_{n+1}\) restrict to \(\rho\). We show that there
is \(b\in G\) for which \((\bar a,b)\) has pattern \(\rho'\).

If the last slot of \(\rho'\) belongs to an existing equality class,
choose any index \(i<n\) in that class and put \(b=a_i\). All such
indices have the same value by \(\rho\). Every equality and edge
involving the last slot is prescribed by that class, so this choice
realizes \(\rho'\).

Otherwise the last slot is a new class. List the distinct entries of
\(\bar a\) once each. Partition this list into \(U,V\) according as
\(\rho'\) requires adjacency or nonadjacency to the last slot. The
extension axiom for these two disjoint lists gives a vertex \(b\)
outside their union, adjacent to every member of \(U\) and to none of
\(V\). It differs from every entry of \(\bar a\); symmetry supplies
the reversed edge instances and irreflexivity supplies its absent loop.
Thus \((\bar a,b)\) has the required pattern. When \(n=0\), both
lists are empty and the \((0,0)\) extension axiom supplies a vertex.
This proves realizability over the \emph{given} tuple, a stronger fact
than realization of the extension pattern somewhere in a model.

Every one-slot extension is of one of these forms. If \(\rho\) has
\(k\) classes there are \(k\) old-class extensions. A new class has
one freely chosen adjacency bit to each old class, so there are
\(2^k\) new-class extensions. The choices are distinct and exhaustive.

\emph{Correctness by induction on the input formula.}
We prove the displayed truth invariant simultaneously for every model,
every tuple and every eligible variable list. An atomic equality is
true exactly when its slots share a class. An atomic edge is true
exactly when its classes are distinct and adjacent. These are precisely
the two atomic selection clauses. Truth and falsity satisfy their
clauses, including for the empty tuple. Negation complements the truth
set within all patterns; conjunction intersects the truth sets for a
common list. The induction hypothesis therefore proves these steps.

Suppose \(\varphi=\exists y\,\psi\), with \(y\) fresh for the
current list. If \(\mathcal G\models\psi(\bar a,b)\), the induction
hypothesis puts the pattern \(\rho'\) of \((\bar a,b)\) in
\(Q_{\bar x,y}(\psi)\). Its restriction is the pattern of \(\bar a\),
which the projection step therefore retains. Conversely, if that
pattern is retained, choose a retained \(\rho'\) restricting to it.
The extension argument above realizes \(\rho'\) as \((\bar a,b)\)
inside this same \(\mathcal G\). The induction hypothesis gives
\(\mathcal G\models\psi(\bar a,b)\), hence the existential formula.
Capture-avoiding renaming preserves satisfaction; it ensures this step
does not change the interpretation of any free slot. Unused slots cause
no exception: the induction is stated for lists containing, rather than
equal to, the free-variable set.

Every tuple satisfies exactly its own pattern formula. The invariant
therefore proves equivalence to the returned disjunction.

\emph{Termination and the decision consequences.}
There are finitely many subformulas, and every recursive call uses a
proper subformula. A variable list has finite length; its partitions,
the graphs on their classes and the resulting pattern formulas form
finite effectively enumerable sets. Complement, intersection and
restriction are finite computable operations on these encodings.
The recursion and the output construction consequently terminate.

If the computed set is nonempty, take one of its patterns. It is
realized in a model of \(T_{\rm RG}\) by
\cref{thm:p03-w03-p03-015-graph-realization}; the invariant makes that
realization satisfy the formula. If the set is empty, the invariant
excludes satisfaction in every model. For validity, equality of the
computed set with \(\Gamma_n\) gives truth at every tuple. If a pattern
is omitted, its realization refutes the formula, so the universal
closure is not a consequence. For \(n=0\) there is exactly one pattern
and each model has the empty tuple. Thus the computed set decides which
sentences hold in every model and which fail in every model. The
procedure's finite enumeration supplies no smaller complexity bound.
\end{proof}
\end{annexitem}
